\section{Finite-width circuit seeds from Gaussian mean shifts}
\label{sec:critical-finite-width-circuit-prefix}

A Boolean circuit can be planted into a finite Gaussian network using
small, prescribed groups of neurons.  Every weight retains its original
Gaussian variance.  Only the means change, and the total relative
entropy of these changes is small.  Saturation corrects the approximation
error at each gate; no exact finite-state repair is needed.

\begin{lemma}[A finite-width Boolean circuit prefix]
\label{lem:critical-finite-circuit-prefix}
Let $H^{(0)}(x)=(1,x_1,\ldots,x_d)$ and let
$H^{(j)}(x)\in\{-1,1\}^{C_j}$, $1\le j\le h$, be a layered
threshold circuit with a constant wire at every level.  Suppose
\[
 H_i^{(j)}(x)=\operatorname{sign}(w_{j,i}\cdot H^{(j-1)}(x)),
 \qquad |w_{j,i}\cdot H^{(j-1)}(x)|\ge1,
 \qquad \|w_{j,i}\|_2\le B,
\]
where $B\ge1$.  Put $C=\max_j C_j$ and $W=\max(d+1,C)$.
For $L\ge h+1$, $0<\delta<1$, and
$0<\eta\le(4B\sqrt W)^{-1}$, define
\[
 a=4+\log(8/\eta),\qquad p_0=1/d,\qquad
 p_j=\frac{p_{j-1}}{8C_jhLa^2B^2},\qquad k_j=\lceil Np_j\rceil.
\]
If
\begin{equation}\label{eq:critical-finite-prefix-width}
 Np_h\ge16\eta^{-2}
       \left[d+\log\frac{8h(C+1)}\delta\right],
\end{equation}
there is a product Gaussian law $P$ on the first $h$ hidden weight
matrices such that
\[
 D(P\Vert\mu_N^{\rm prefix})\le\frac{N}{8L}.
\]
Here the covariance of every weight row is unchanged from the critical
prior.  Under $P$, with probability at least $1-\delta$, the following
hold simultaneously on every input in $\{-1,1\}^d$.
Each circuit wire $i$ at layer $j$ has a prescribed group of $k_j$
neurons whose empirical mean $m_{j,i}$ satisfies
\[
 \|m_{j,i}-H_i^{(j)}\|_\infty\le\eta.
\]
If $q_j(x)=N^{-1}\|h_j(x)\|_2^2$, set
\begin{align*}
 \theta_j&=C_jk_j/N, & \bar q_0&=(d+1)/d,\\
 \bar q_j&=(1-\theta_j)F(\bar q_{j-1})+\theta_j,
 & F(q)&=\mathbb E\tanh^2(\sqrt qZ).
\end{align*}
Then $\|q_j-\bar q_j\|_\infty\le2j\eta$.
If a final wire is $f$, there is a deterministic unit vector
$u\in\mathbb R^N$ such that, writing $p=k_h/N$ and
$c(x)=N^{-1/2}\langle u,h_h(x)\rangle$,
\begin{align}
 |c(x)-\sqrt p f(x)|&\le\sqrt p\eta,
 \label{eq:critical-finite-prefix-score}\\
 |q_h(x)-c(x)^2-(\bar q_h-p)|&\le(2h+3)\eta,
 \qquad \bar q_h-p\ge0.
 \label{eq:critical-finite-prefix-residual}
\end{align}
Consequently its scalar-channel metric error relative to the ideal seed
$(v,r)=(p,\bar q_h-p)$ is at most $\sqrt{(2h+4)\eta}$.
\end{lemma}

\begin{proof}
Choose disjoint deterministic groups of size $k_j$, one for each wire
at layer $j$.  They fit within the width: the hypothesis gives
$Np_j\ge1$, and hence
\[
 \theta_j\le2C_jp_j\le\frac1{4hLa^2B^2}\le\frac14.
\]
Rows outside these groups retain their original Gaussian law.  A row
in the first-layer group for gate $i$ has mean $aw_{1,i}$ and covariance
$I/d$.  At layer $j\ge2$, the row mean is $aw_{j,i,g}/k_{j-1}$ on
each coordinate of the preceding group for wire $g$, and zero on the
remaining coordinates; its covariance is $I/N$.  These weight means
are deterministic, so all the weight rows remain independent under $P$.
The resulting preactivation mean is exactly
\[
 a\sum_g w_{j,i,g}m_{j-1,g}(x).
\]
The relative entropy of each controlled row is at most
$a^2B^2/(2p_{j-1})$, including the first layer with $p_0=1/d$.
For later layers this follows directly from
\[
 \frac N2\sum_g k_{j-1}
       \left(\frac{aw_{j,i,g}}{k_{j-1}}\right)^2
 \le\frac{a^2B^2}{2p_{j-1}}.
\]
Gaussian additivity therefore gives
\[
 \frac1N D(P\Vert\mu_N^{\rm prefix})
 \le\sum_{j=1}^h\frac{\theta_ja^2B^2}{2p_{j-1}}
 \le\sum_{j=1}^h\frac1{8hL}=\frac1{8L}.
\]
This bound does not depend on any concentration event.

Suppose the preceding empirical gate means have error at most $\eta$.
Every controlled preactivation then has the correct sign and mean
magnitude at least
\[
 a(1-B\sqrt W\eta)\ge3a/4\ge a/2.
\]
Its variance is at most two in the first layer and at most one
thereafter.  For a Gaussian variable $V$ of variance at most two whose
mean has absolute value at least $a/2$, Gaussian exponential moments
and $1-\tanh t\le2e^{-2t}$ give
\[
 |\mathbb E\tanh V-\operatorname{sign}(\mathbb EV)|
       \le2e^{-a+4}\le\eta/4,
 \qquad
 1-\mathbb E\tanh^2V\le4e^{-a+4}\le\eta/2.
\]
Thus each new gate mean is close to its Boolean value independently
of the size of the preceding error, provided that error stays below
the stated gate-margin threshold.

Conditionally on the preceding layers, the rows in each group are
independent and identically distributed.  Hoeffding's inequality,
followed by a union over all $2^d$ inputs, bounds the chance that any
empirical gate mean differs from its conditional mean by more than
$\eta/2$ by
\[
 2C_j2^d\exp(-k_j\eta^2/8).
\]
The resulting error relative to the Boolean wire is less than $\eta$.
The squared features are independent, bounded in $[0,1]$, and need
not have identical laws.  Their empirical average differs from its
conditional mean by at most $\eta$ everywhere except with probability
at most $2\cdot2^d\exp(-2N\eta^2)$.

The function $F$ is one-Lipschitz on $[0,\infty)$: its heat derivative
is
\[
 F'(q)=\mathbb E\bigl[1-4\tanh^2(\sqrt qZ)
                              +3\tanh^4(\sqrt qZ)\bigr],
\]
whose absolute value is at most one.  The conditional diagonal mean
is within $\theta_j\eta/2$ of
$(1-\theta_j)F(q_{j-1}(x))+\theta_j$.  Therefore, on the
concentration events,
\[
 |q_j(x)-\bar q_j|
 \le |q_{j-1}(x)-\bar q_{j-1}|+3\eta/2\le2j\eta.
\]
Conditional union bounds over the layers and
\eqref{eq:critical-finite-prefix-width} prove that all these events
hold with probability at least $1-\delta$.

Take $u$ to be the target-group indicator divided by $\sqrt{k_h}$.
Then $\|u\|_2=1$ and $c=\sqrt p\,m_{h,\mathrm{target}}$, giving
\eqref{eq:critical-finite-prefix-score}.  Also
$|c^2-p|\le3p\eta$, which proves
\eqref{eq:critical-finite-prefix-residual}; its deterministic residual
is nonnegative because $\bar q_h\ge\theta_h=C_hp\ge p$.
Finally $(\sqrt u-\sqrt v)^2\le|u-v|$ gives the stated metric
bound.  Orthogonally projecting the next Gaussian weight row onto $u$
yields an exact independent Gaussian score/residual decomposition,
not an approximate positive-semidefinite covariance inequality.
\end{proof}

\begin{remark}[Seed-relative accuracy and the circuit depth parameter]
\label{rem:critical-finite-prefix-accuracy}
Put
\[
 R=d+C+B+h+L+2,\qquad
 T=\log(d+2)+h\log(C+B+h+L+2).
\]
If subsequent scalar steering needs $\eta\le cR^{-b}p_h^b$ for
some fixed universal $b$, choose $a=KT$ and $\eta=8e^{4-a}$ with
a sufficiently large universal $K$.  The gate-margin condition then
holds.  Moreover, since $T\ge h\log(h+2)$,
\[
 h\log T=h\log h+h\log(T/h)\le2T.
\]
Consequently
\[
 \log(1/p_h)=\log d+
       \sum_{j=1}^h\log(8C_jhLa^2B^2)
       \le C_0(1+\log K)T.
\]
The required accuracy inequality follows by taking $K$ sufficiently
large: the coefficient in $\log(1/\eta)$ is linear in $K$, whereas
the coefficient on the right is logarithmic in $K$.  The same reasoning
allows any fixed requirement $\eta\le e^{-c_1T}$.
For $\delta^{-1}\le e^{c_2T}$ with a fixed universal $c_2$, the
width condition \eqref{eq:critical-finite-prefix-width} is therefore
satisfied by $N\ge e^{C_1T}$, for a universal constant $C_1$ depending
on the required fixed accuracy exponents.  In particular this is
$R^{O(h)}$.  Arbitrary smaller failure probabilities contribute the
additional factor $1+\log(1/\delta)$ in the explicit width condition.
\end{remark}
